% Vista editorial exacta de corpus/source/masas/06_operadores.tex, líneas 50--113. Las líneas 1--29 y 30--49 ya residen exactamente en 73_rev2_electron_lectores_torres.tex.
\subsection{Lectores bidireccionales \(K_D\) y \(K_\Omega\): perfiles, coincidencias y operadores por multisección}
La ruta orientada produce dos familias de invariantes. Para los defectos temporales \(D_j\) y la memoria de retorno \(\Omega\), definimos

\[
K_D=\left(D_{27}+D_{36},D_1,\frac{D_4-D_3}{2}\right),
\qquad
K_\Omega=(\Omega,54,P_6).
\]
Sobre las 324 rutas aparecen 14 perfiles distintos de \(K_D\), 9 perfiles de \(K_\Omega\), 40 pares conjuntos y 18 coincidencias. El único perfil común de la ruta electrónica es

\[
K_D(\Gamma_e)=K_\Omega(\Gamma_e)=(80,54,6).
\]
Estos lectores no son etiquetas decorativas. \(K_D\) mide discrepancias de la palabra; \(K_\Omega\) mide el retorno enriquecido. Su igualdad en el centro certifica la clausura bidireccional de rango uno. Fuera del centro producen operadores de fibra cuyos espectros completos deben publicarse aun cuando una realización física posterior seleccione una sección escalar.

\Needspace{7\baselineskip}
\subsection{Refinamiento armónico propio de cada estado}
El refinamiento completo vive en el espacio de coeficientes

\[
\mathcal E_{90,120}
=\left\{\eta=(\eta_h)_{h\in\langle90,120\rangle}:
\sum_h|\eta_hs_h|<\infty\right\}.
\]
La aplicación buscada no recibe una masa. Recibe la ruta enriquecida:

\[
\mathfrak T:\Gamma_{\rm enr}\longrightarrow\mathcal E_{90,120},
\qquad
\Gamma\longmapsto\eta(\Gamma).
\]
Sus coordenadas deben ser funciones explícitas de retornos, acarreo, devanado, memoria y frontera. La condición de naturalidad es

\[
\eta(r_{n,m}\Gamma)=r^{\mathcal E}_{n,m}\eta(\Gamma),
\]
y la convergencia debe acompañarse de una cota de cola calculable. El carácter completo actúa entonces por

\[
\widehat M_F
=\sum_{\gamma\in F}
m_{\rm clk}\,
\chi_{\rm full}(\sigma(\gamma),\eta(\gamma))
|\gamma\rangle\langle\gamma|.
\]
La capacidad de reconstruir una diferencia numérica dada demuestra completitud representativa de la jerarquía; no demuestra por sí sola que \(\mathfrak T\) haya elegido sus coeficientes. La diferencia entre ambas afirmaciones queda preservada en cada fila cuantitativa.

\Needspace{7\baselineskip}
\subsection{Escalarización por acoplamiento}
Sea \(P_{X,a}\) el proyector determinado por el estado \(X\) y el aparato o canal físico \(a\). La distribución de masa es

\[
\mu_{X,a}(B)
=\langle\Psi_X,
P_{X,a}E_{\widehat M_F}(B)P_{X,a}\Psi_X\rangle.
\]
Hay una línea escalar cuando la compresión satisface

\[
P_{X,a}\widehat M_FP_{X,a}=m_XP_{X,a}.
\]
Éste es el significado preciso de «medir es acoplarse». La cifra no estaba oculta en una celda nominal; aparece como autovalor de una compresión física de un operador ya construido. Si la compresión conserva varios autovalores, la predicción correcta es un multiplete o una distribución, no una media elegida para coincidir con una referencia.

\Needspace{7\baselineskip}

